

%% 
%% This is a LaTeX document generated by automatic conversion of a
%% Mathematica notebook using Mathematica 4.0.
%% 
%% This document uses special macros that are defined in a LaTeX 
%% package file with a name of the form notebook*.sty.  Appropriate 
%% commands for loading the package have been included in this document.
%% The LaTeX package files can be found in a directory with path:
%% 
%% /usr/local/mathematica/SystemFiles/IncludeFiles/TeX/
%% 
%% The LaTeX package used in this document may require that your 
%% LaTeX implementation support fonts other than the standard Computer 
%% Modern fonts that are used by LaTeX.  See the Additional Information 
%% section under the TeXSave entry in the Mathematica Reference Guide 
%% for more information.
%% 
%% Implementation-specific instructions for installing TeX-related files
%% can be found at this URL.
%% 
%% http://www.wolfram.com/support/FrontEnds/Export/TeX
%% 
%% 

\documentclass{article}
\usepackage{notebook2e, latexsym}


\begin{document}

\Title{FOURIER SERIES}
\Subtitle{LECTURE 9 }
\Subsubtitle{ELE314: LINEAR SIGNALS AND SYSTEMS}
\Subsubtitle{Mariusz Jankowski \copyright{}

University of Southern Maine

mjkcc@usm.maine.edu}
This loads some useful packages.

\dispSFinmath{<<\Mvariable{SignalPlot`}}
\dispSFinmath{<<\Mvariable{Graphics`Graphics`}}
\begin{CellGroup}
\Section*{INTRODUCTION}
Consider the following important fact about LTI systems: the system response to a linear combination of some basic signals is a superposition of
the individual responses of the system to each basic signal. Furthermore, the response of an LTI system to a sinusoidal signal is sinusoidal in form!
Formally, the family of sinusoids is an eigenfunction of linear time invariant systems. Sinusoidal signals are therefore extremely useful in the
analysis of LTI systems. In the early 19th century, the French mathematician Jean Baptiste Fourier showed that a large class of interesting functions
could be represented by linear combinations of harmonically related sinusoids. 

\mathGraphic{lecture09.tex_gr1.eps}
The complex Fourier series is a representation of a periodic signal in terms of linear combinations of harmonically related complex exponentials:

	\inlineTFinmath{x(t)\multsp =\multsp \sum _{k=-\infty }^{\infty }{c_k}\multsp \multsp {{\ExponentialE }^{\multsp \multsp j\multsp \frac{2\pi }{T}k\multsp
t}}}

where T is the period of the signal \inlineTFinmath{x(t)}. The quantity \inlineTFinmath{\frac{2\pi }{T}=\multsp 2\pi \multsp {f_T}\multsp ={{\omega
}_T}\multsp } is called the fundamental frequency (in radians/sec). All other complex exponentialss have frequencies that are integer multiples of
the fundamental. The coefficients \inlineTFinmath{{c_k}} are known as the {\itshape Fourier series coefficients} of the signal \inlineTFinmath{x(t)}
and are generally complex. Equation (1) is called the {\itshape synthesis equation}, since given a set of FS coefficients (possibly infinite) it
provides a formula for "synthesizing" the continuous time signal \inlineTFinmath{x(t)}. Immediately, two important questions arise:

(1) What kind of signals posses a FS expansion?

(2) How to obtain the FS coefficients \inlineTFinmath{{c_k}}?

You will be asked to consider the first of the questions in a homework problem. Here we turn to the second of these questions. The "inverse" relationship
between a continuous-time signal and its FS coefficients is not very difficult to derive and the specifics can be found in most introductory texts
on the subject [Oppenheim, pp. 190-191]. Here we simply state the {\itshape analysis formula}:



\inlineTFinmath{{c_k}=\multsp \frac{1}{T}\int _{-\frac{T}{2}}^{\frac{T}{2}}\multsp x(t)\multsp \multsp {{\ExponentialE }^{-\multsp j2\frac{\pi \multsp
}{T}\multsp k\multsp t}}\DifferentialD t}

Equations (1) and (2) define the so-called {\itshape complex} {\itshape Fourier series} formulation. Other related formulations found in the literature
are the trigonometric series formulation and the cosine series formulation [Kreyszig, 1972].



\begin{CellGroup}
\Exercise{Problem 1}
\ExerciseText{Using Euler's identity, \inlineTFinmath{{{\ExponentialE }^{\multsp \ImaginaryI \multsp \theta }}\multsp =\multsp \multsp \cos(\theta
)\multsp +\multsp \ImaginaryI \multsp \sin(\theta )}, rewrite (1) in terms of sine and cosine functions.}



\end{CellGroup}
We now turn our attention to the problem of evaluating (1) given a specific periodic signal. Assume signal \inlineTFinmath{x(t)} is a pulse train,
with a period \inlineTFinmath{T=1}, and pulse width  \(\tau \)\(=\) \inlineTFinmath{\frac{T}{2}}, centered on \inlineTFinmath{t=0}.



\dispSFinmath{\Muserfunction{fun}[\Mvariable{t\_}]:=\Mfunction{UnitStep}[t]-\Mfunction{UnitStep}[t-0.5]}
\begin{CellGroup}
\dispSFinmath{\MathBegin{MathArray}{l}
\Mfunction{Plot}[\Muserfunction{fun}[\Mfunction{Mod}[t+1/4,1]],\{t,-2,2\},  \\
\noalign{\vspace{0.5ex}}
\hspace{1.em} \Mvariable{PlotRange}\rightarrow \Mvariable{All},\Mvariable{PlotLabel}\rightarrow \Mvariable{Tren\multsp de\multsp Pulsos}]\\
\MathEnd{MathArray}}
\mathGraphic{lecture09.tex_gr2.eps}
\dispSFoutmath{\SkeletonIndicator \Mvariable{Graphics}\SkeletonIndicator }

\end{CellGroup}
The FS coefficients are evaluated using (2).

\begin{CellGroup}
\dispSFinmath{{c_{\Mvariable{k\_}}}\multsp =\int _{-\multsp \frac{1}{4}}^{\frac{1}{4}}\multsp {{\ExponentialE }^{\multsp -\ImaginaryI \multsp 2\multsp
\pi \multsp k\multsp t}}\DifferentialD t\multsp //\Mfunction{ExpToTrig}}
\dispSFoutmath{\frac{\sin \big[\frac{k\multsp \pi }{2}\big]}{k\multsp \pi }}

\end{CellGroup}
For the special case of \inlineTFinmath{k=0}, we get \inlineTFinmath{{c_0}=\multsp \frac{1}{T}\multsp {{\int }_T}\multsp x(t)\DifferentialD t\multsp
=\multsp \frac{1}{2}}, or equivalently as the limit of the expression defining \inlineTFinmath{{c_k}}, as k tends to zero. 



\begin{CellGroup}
\dispSFinmath{{c_0}=\Mfunction{Limit}[{c_k},k->0]}
\dispSFoutmath{\frac{1}{2}}

\end{CellGroup}
Here are a few of the values of the FS coefficients of the periodic pulse train.

\begin{CellGroup}
\dispSFinmath{\MathBegin{MathArray}{l}
\Muserfunction{LinePlot}[\Mfunction{Table}[\{k,{c_k}\},\{k,-21,21\}],\Mvariable{PlotRange}->\Mvariable{All},  \\
\noalign{\vspace{0.5ex}}
\hspace{1.em} \Mvariable{PlotLabel}\rightarrow \Mvariable{Coeficientes\multsp serie\multsp de\multsp Fourier\multsp de\multsp tren\multsp de\multsp
pulsos}];\\
\MathEnd{MathArray}}
\mathGraphic{lecture09.tex_gr3.eps}

\end{CellGroup}
This graphical representation of the Fourier series coefficients is called a line spectrum. Typically the magnitude is plotted as a function of the
integer index \inlineTFinmath{k}. The frequency corresponding to \inlineTFinmath{k=1} is called the fundamental harmonic, which in the present case
is \inlineTFinmath{\frac{2\pi }{T}=\frac{2\pi }{1}} \inlineTFinmath{\frac{\Mvariable{radians}}{\Mvariable{sec}}} or 1 Hz. Note the symmetry in the
values of the coefficients and the fact that they are REAL. Both characteristics are a direct consequence of the properties of the original time-domain
signal.



\begin{CellGroup}
\Exercise{Problem 2}
\ExerciseText{Recalculate the Fourier series coefficients of the signal \inlineTFinmath{x(t)}, keeping the period fixed and decreasing the pulse
width to \inlineTFinmath{\tau =\frac{1}{8}} and \inlineTFinmath{\frac{1}{64}}. Discuss the effect on the FS line spectrum.}



\end{CellGroup}

\end{CellGroup}
\begin{CellGroup}
\Section*{SIGNAL RECONSTRUCTION}
Any periodic function that satisfies the so-called {\itshape Dirichlet conditions} has a Fourier series representation in terms of linear combinations
of harmonically related sinusoids. As the synthesis formula (2) shows, a periodic signal may be an infinite sum of harmonically related sinusoids.
For practical reasons, it is of interest to investigate the approximation of a periodic signal with a FINITE sum.



\inlineTFinmath{{x_M}(t)\multsp =\multsp \sum _{k=-M}^{M}\multsp {c_k}\multsp \multsp {{\ExponentialE }^{\multsp j\multsp {{\omega }_T}\multsp k\multsp
t}}}

where \inlineTFinmath{{x_M}} is the approximating signal, synthesized from \inlineTFinmath{2M+1} FS coefficients. Given the FS coefficients calculated
in the earlier example, we obtain approximations of the original signal using \inlineTFinmath{M=1,5,11,21} coefficients.



\dispSFinmath{{x_{\Mvariable{M\_}}}[\Mvariable{t\_}]:=\sum _{k=-M}^{M}{c_k}\multsp {{\ExponentialE }^{\multsp \ImaginaryI \multsp 2\pi \multsp k\multsp
t}};}
\begin{CellGroup}
\dispSFinmath{\MathBegin{MathArray}{l}
\Muserfunction{DisplayTogetherArray}[\{\{\Mfunction{Plot}[{c_0},\{t,-1,1\},\Mvariable{PlotLabel}\rightarrow \Mvariable{continua}],  \\
\noalign{\vspace{0.5ex}}
\hspace{3.em} \Mfunction{Plot}[{x_4}[t],\{t,-1,1\},\Mvariable{PlotLabel}\rightarrow 4\multsp armonicas\multsp y\multsp continua]\},  \\
\noalign{\vspace{0.5ex}}
\hspace{2.em} \{\Mfunction{Plot}[{x_{16}}[t],\{t,-1,1\},\Mvariable{PlotLabel}\rightarrow 16\multsp armonicas\multsp y\multsp continua],  \\
\noalign{\vspace{0.5ex}}
\hspace{3.em} \Mfunction{Plot}[{x_{32}}[t],\{t,-1,1\},\Mvariable{PlotLabel}\rightarrow 32\multsp armonicas\multsp y\multsp continua]\}\}];\\
\MathEnd{MathArray}}
\mathGraphic{lecture09.tex_gr4.eps}

\end{CellGroup}
Clearly, as the number of coefficients increases, the approximation improves. The interesting question to ask is whether this result holds true for
any periodic function? Also, how good is the approximation in some general sense? A mathematically rigorous treatment of Fourier theory provides
answers to these important questions. It is well known that the FS yields the best approximation in the mean squared error sense, to periodic signals
that satisfy the Dirichlet conditions. In addition, for continuous periodic signals, the approximating signal converges to the original (zero error
at every value of t), while for signals with discontinuities, the approximating signal converges to the average value at the discontinuity and to
the signal everywhere else. For finite sums, the behaviour of the approximating signal exhibits ripples and overshoots at a discontinuity. This is
known as the {\itshape Gibb's} phenomenon.  In order to determine how good a particular approximation is, we need to specify a quantitative measure
of the size of the approximation error. A commonly used criterion is the energy over one period or average power of the error signal:



\inlineTFinmath{{e_M}\multsp =\multsp \frac{1}{T}\int _{-\frac{T}{2}}^{\frac{T}{2}}|{e_M}(t){|^2}\DifferentialD t}

where \inlineTFinmath{{e_M}(t)} denotes the approximation error, defined in the following way



\inlineTFinmath{{e_M}(t)\multsp =\multsp x(t)\multsp -\multsp {x_M}(t)\multsp =\multsp x(t)\multsp -\multsp \sum _{k=-M}^{M}{c_k}\multsp {{\ExponentialE
}^{\multsp \multsp j\multsp \multsp {{\omega }_T}\multsp k\multsp t}}}

\begin{CellGroup}
\Exercise{Problem 3}
\ExerciseText{Plot the reconstruction error signal  \inlineTFinmath{{e_M}(t)\multsp =\multsp x(t)-{x_M}(t)}, for  \inlineTFinmath{M=11,21} on the
interval \inlineTFinmath{0\leq t\leq 0.5}.}



\end{CellGroup}

\end{CellGroup}
\begin{CellGroup}
\Section*{SIGNAL POWER}
Periodic, continuous-time signals are also known as power signals. The following relation, known as {\itshape Parseval's relation} for periodic signals,
establishes an equality between time and frequency-domain evaluation of average signal power. 



\inlineTFinmath{{P_T}=\frac{1}{T}\int _{-\frac{T}{2}}^{\frac{T}{2}}|x(t){|^2}\DifferentialD t\multsp =\multsp \sum _{k=-\infty }^{\infty }|{c_k}\multsp
{|^2}}	

The derivation may be found in Oppenheim, p. 205. As an example, consider using (6) to calculate average power in the periodic pulse train. The time
domain calculation is trivial and by inspection the signal power is \inlineTFinmath{{P_T}}\(=\)\inlineTFinmath{\tau }. However, the rhs of (6) requires
an infinite summation. Consider therefore an approximation to the rhs by evaluating the sum over a finite number of coefficients. Intuitively, as
the number of coefficients is increased, since \inlineTFinmath{{x_M}(t)\multsp \rightarrow \multsp x(t)}, then the sum should converge to \inlineTFinmath{{P_T}}.
Here we verify that this indeed happens. 



\dispSFinmath{\MathBegin{MathArray}{l}
{P_T}\multsp =\multsp 0.5;  \\
\noalign{\vspace{1.0625ex}}  \\
{P_{\Mvariable{M\_}}}=\sum _{k=-M}^{M}{{\Mfunction{Abs}[{c_k}]}^2};
\MathEnd{MathArray}}
where \inlineTFinmath{{P_M}} is the average power of the approximation signal \inlineTFinmath{{x_M}(t)}. Here are a few values of the average power
as a function of \inlineTFinmath{M}.



\begin{CellGroup}
\dispSFinmath{\Mfunction{Table}[{P_M},\multsp \{M,11\}]//\Mfunction{N}}
\dispSFoutmath{\MathBegin{MathArray}{l}
\{0.452642,0.452642,0.475158,0.475158,0.483264,  \\
\noalign{\vspace{0.5ex}}
\hspace{1.em} 0.483264,0.487399,0.487399,0.489901,0.489901,0.491576\}\\
\MathEnd{MathArray}}

\end{CellGroup}
The ratio of signal power in the approximation and the original signal defines a common measure of the quality of the approximation. The closer this
number is to unity the better is the FS representation of the original signal (in the sense of mean squared error). 

\begin{CellGroup}
\dispSFinmath{\frac{\%}{{P_T}}}
\dispSFoutmath{\MathBegin{MathArray}{l}
\{0.905285,0.905285,0.950316,0.950316,0.966528,  \\
\noalign{\vspace{0.5ex}}
\hspace{1.em} 0.966528,0.974799,0.974799,0.979802,0.979802,0.983152\}\\
\MathEnd{MathArray}}

\end{CellGroup}
The approximation error average power \inlineTFinmath{{e_M}} defined in (4) is clearly the difference between total power and power in the approximating
signal. We use Parseval's relation to calculate \inlineTFinmath{{e_M}} for selected values of \inlineTFinmath{M}.



\begin{CellGroup}
\dispSFinmath{\MathBegin{MathArray}{l}
{e_{\Mvariable{M\_}}}:={P_T}-{P_M};  \\
\noalign{\vspace{0.5ex}}  \\
\Mfunction{Table}[{e_M},\multsp \{M,21\}]//\Mfunction{N}
\MathEnd{MathArray}}
\dispSFoutmath{\MathBegin{MathArray}{l}
\{0.0473576,0.0473576,0.0248418,0.0248418,0.0167361,0.0167361,  \\
\noalign{\vspace{0.5ex}}
\hspace{1.em} 0.0126006,0.0126006,0.0100988,0.0100988,0.00842407,  \\
\noalign{\vspace{0.5ex}}
\hspace{1.em} 0.00842407,0.00722501,0.00722501,0.00632437,0.00632437,  \\
\noalign{\vspace{0.5ex}}
\hspace{1.em} 0.00562319,0.00562319,0.00506185,0.00506185,0.00460235\}\\
\MathEnd{MathArray}}

\end{CellGroup}
Here is the avearge power of the error signal as a function of \inlineTFinmath{M}, 



\begin{CellGroup}
\dispSFinmath{\Muserfunction{LinePlot}[\%,\multsp \Mvariable{PlotRange}->\{0,\multsp 0.1\}];}
\mathGraphic{lecture09.tex_gr5.eps}

\end{CellGroup}
\begin{CellGroup}
\Exercise{Problem 4}
\ExerciseText{Obtain \inlineTFinmath{{e_5}} using (4) and (5) directly.}



\end{CellGroup}

\end{CellGroup}
\begin{CellGroup}
\Section*{PROBLEM SOLUTIONS}
\begin{CellGroup}
\Exercise{Problem 1}
\ExerciseText{Using Euler's identity, \inlineTFinmath{{{\ExponentialE }^{\multsp \ImaginaryI \multsp \theta }}\multsp =\multsp \multsp \cos(\theta
)\multsp +\multsp \ImaginaryI \multsp \sin(\theta )}, rewrite (1) in terms of sine and cosine functions.}



\end{CellGroup}
\begin{CellGroup}
\Exercise{Problem 2}
\ExerciseText{Recalculate the Fourier series coefficients of the signal \inlineTFinmath{x(t)}, keeping the period fixed and decreasing the pulse
width to \inlineTFinmath{\tau =\frac{1}{16}} and \inlineTFinmath{\frac{1}{64}}. Discuss the effect on the FS line spectrum.}



\end{CellGroup}
\ExerciseText{{\bfseries Solution}. We scale each signal (and therefore each coefficient) to keep the average value constant and equal to 1.}


\dispSFinmath{\MathBegin{MathArray}{l}
\Muserfunction{c1}[\Mvariable{k\_}]:=8\frac{\sin \big[\frac{\multsp k\multsp \pi }{8}\big]}{k\multsp \pi };\multsp \Muserfunction{c2}[\Mvariable{k\_}]:=32\frac{\sin
\big[\frac{k\multsp \pi }{32}\big]}{k\multsp \pi };  \\
\noalign{\vspace{1.07292ex}}  \\
\Muserfunction{c1}[0]=1;\multsp \Muserfunction{c2}[0]:=1;
\MathEnd{MathArray}}
\begin{CellGroup}
\dispSFinmath{\MathBegin{MathArray}{l}
\Mvariable{DisplayTogetherArray}[  \\
\noalign{\vspace{0.5ex}}
\hspace{1.em} \{\Muserfunction{LinePlot}[\Mfunction{Table}[\{k,\Mfunction{Abs}[\Muserfunction{c1}[k]]\},\multsp \{k,0,63\}],\multsp \Mvariable{PlotRange}\rightarrow
\Mvariable{All}],  \\
\noalign{\vspace{0.5ex}}
\hspace{2.em} \Muserfunction{LinePlot}[\Mfunction{Table}[\{k,\Mfunction{Abs}[\Muserfunction{c2}[k]]\},\multsp \{k,0,63\}],\multsp \Mvariable{PlotRange}\rightarrow
\Mvariable{All}]\}];\\
\MathEnd{MathArray}}
\mathGraphic{lecture09.tex_gr6.eps}

\end{CellGroup}
For narrower pulses, the FS coefficient spectrum "widens". The position of the first zero moves to a higher harmonic, and thus the "mainlobe" widens.

\begin{CellGroup}
\Exercise{Problem 3}
\ExerciseText{Plot the reconstruction error signal  \inlineTFinmath{{e_M}(t)\multsp =\multsp x(t)-{x_M}(t)}, for  \inlineTFinmath{M=11,21}.}


\ExerciseText{{\bfseries Solution}. The signal has a value of one for \inlineTFinmath{0\leq t\leq 0.25} and zero for the remainder of the region
of interest. For plotting purposes I define \inlineTFinmath{x(t)} as}


\dispSFinmath{\MathBegin{MathArray}{l}
x[\Mvariable{t\_}]:=0;  \\
\noalign{\vspace{0.5ex}}  \\
x[\Mvariable{t\_}]:=1\multsp /;\multsp t<0.25
\MathEnd{MathArray}}
\ExerciseText{Here we now plot the approximation error signal.}
\begin{CellGroup}
\dispSFinmath{\Mfunction{Plot}[x[t]-{x_8}[t],\multsp \{t,0,0.5\},\multsp \Mvariable{PlotRange}->\Mvariable{All}];}
\mathGraphic{lecture09.tex_gr7.eps}

\end{CellGroup}
\begin{CellGroup}
\dispSFinmath{\Mfunction{Plot}[x[t]-{x_{16}}[t],\multsp \{t,0,0.5\},\multsp \Mvariable{PlotRange}->\Mvariable{All}];}
\mathGraphic{lecture09.tex_gr8.eps}

\end{CellGroup}

\end{CellGroup}
\begin{CellGroup}
\Exercise{Problem 4}
\ExerciseText{Obtain \inlineTFinmath{{e_5}} using (4) and (5) directly.}



\end{CellGroup}
\ExerciseText{{\bfseries Solution:} Due to symmetry we only need to integrate over the half interval \inlineTFinmath{0\leq t\leq 0.5}. First obtain
the approximation signal,}


\begin{CellGroup}
\dispSFinmath{y={x_4}[t]\multsp //\Mfunction{ExpToTrig}//\Mfunction{Chop}}
\dispSFoutmath{\frac{1}{2}+\frac{2\multsp \cos [2\multsp \pi \multsp t]}{\pi }-\frac{2\multsp \cos [6\multsp \pi \multsp t]}{3\multsp \pi }}

\end{CellGroup}
\ExerciseText{Substitute the approximation signal into (4) and evaluate}
\begin{CellGroup}
\dispSFinmath{\multsp {{\ScriptCapitalE }_4}\multsp =\multsp 2\int _{0.25}^{0.5}{y^2}\DifferentialD t\multsp +\multsp 2\int _{0}^{0.25}{{(1.-y)}^2}\DifferentialD
t}
\dispSFoutmath{0.0248418}

\end{CellGroup}

\end{CellGroup}
\begin{CellGroup}
\Section*{HOMEWORK PROBLEMS}
\begin{CellGroup}
\Exercise{Homework Problem 9.1}
\ExerciseText{One of the fundamental questions of Fourier theory asks which periodic functions posses a FS expansion. Consult your textbook and references
for the correct answer.}

\end{CellGroup}
\begin{CellGroup}
\Exercise{Homework Problem 9.2}
\ExerciseText{Obtain the FS coefficients of a periodic pulse train, but one that is shifted to the right, so that the pulse begins at \inlineTFinmath{t=0}.
Comment on any similarities and differences between the present result and the one discussed in class.}



\end{CellGroup}
\begin{CellGroup}
\Exercise{Homework Problem 9.3}
\ExerciseText{Determine the Fourier series coefficients for each of the following periodic signals:

(a) see Figure P3.22b, Oppenheim/Willsky, p. 255

(b) see Figure P3.22e, Oppenheim/Willsky, p. 256}

\end{CellGroup}
\begin{CellGroup}
\Exercise{Homework Problem 9.4}
\ExerciseText{Estimate the magnitude of the overshoot in a finite series approximation of an ideal pulse. Does this magnitude depend on the number
of coefficients used?}

\end{CellGroup}
\begin{CellGroup}
\Exercise{Homework Problem 9.5}
\ExerciseText{Compare the Fourier series of two periodic signals by plotting the approximation error as a function of number of FS coefficients \inlineTFinmath{M}.
Vary \inlineTFinmath{M} from 1 to 5.

(a) pulse train used in class example,

(b) pulse train with period T\(=\)1, formed from exponentials of the form}


\dispSFinmath{g[\Mvariable{t\_}]:=\exp [-\multsp 25\multsp {t^2}]}

\end{CellGroup}

\end{CellGroup}

\end{document}
